\section{Exact identification does not imply a power-law stability bound}
\label{sec:shell-instability}
The scalar two-circle uniqueness theorem is prior work of Perez
\cite{Perez2021}. We now distinguish that exact theorem, and the preceding
finite-rank extension, from quantitative recovery with errors. All norms below
are on a fixed disk; the measurement circles and the target point are fixed.

\begin{theorem}[A two-circle instability sequence]\label{thm:shell-instability}
For every integer $M\ge1$ there are scalar polynomials $p_M,q_M$ of degree
at most $8M+2$ such that
\[
 \tfrac12\le\|p_M\|_{H^\infty(D_2)}<1,\qquad
 \tfrac12\le\|q_M\|_{H^\infty(D_2)}<1,
\]
and, for $H_M=|p_M|^2-|q_M|^2$,
\begin{align}
 H_M(z)&=0 &&(|z|=1),\\
 \sup_{|z|=1/2}|H_M(z)|&\le 2^{-4M^2-18M-3},\\
 |H_M(3/4)|&\ge\frac1{40}(3/8)^{8M+2}.
\end{align}
In particular there are no constants $C,\alpha>0$ such that
\[
 |\,|p(3/4)|^2-|q(3/4)|^2\,|
 \le C\left(\max_{\rho\in\{1/2,1\}}\sup_\theta
 |\,|p(\rho e^{i\theta})|^2-|q(\rho e^{i\theta})|^2\,|\right)^\alpha
\]
for all scalar holomorphic $p,q$ bounded by one on $D_2$. This fails even
when both functions are polynomials and each has disk supremum at least $1/2$.
\end{theorem}
\begin{proof}
Put $b=1/4$, $L=4M+1$, and define the Laurent polynomials
\[
 \begin{gathered}
 A_M(z)=\sum_{n=-M}^M b^{2n^2+n}z^{4n+1},\qquad B_M(z)=A_M(1/z),\\
 p_M(z)=\frac{z^L A_M(z)}{4\,2^L},\qquad
 q_M(z)=\frac{z^L B_M(z)}{4\,2^L}.
 \end{gathered}
\]
The latter two are ordinary polynomials of degree at most $8M+2$.
All coefficients are positive. At $|z|=2$ the Laurent sums are bounded by
\[
 A_\infty(2)=2\sum_{n\in\mathbb Z}2^{-4n^2+2n}
 \le2\left(1+2\sum_{n\ge1}2^{-2n}\right)=\frac{10}3<4.
\]
Changing $n$ to $-n-1$ gives $A_\infty(1/2)=A_\infty(2)$, so this bound
also applies to $B_M(2)$. The maximum principle proves the disk upper bound.
The $n=0$ term of $A_M(2)$ is $2$ and the $n=-1$ term of $A_M(1/2)$ is
$2$ for $M\ge1$, proving the lower bound for both disk suprema.

Reindexing the finite sum leaves exactly two boundary terms:
\begin{equation}\label{eq:theta-truncation-defect}
 D_M(z):=A_M(bz)-b^{-2}z^{-4}A_M(z)
 =b^{2M^2+5M+1}z^{4M+1}-b^{2M^2-M-2}z^{-4M-3}.
\end{equation}
Consequently, on $|z|=2$,
\[
 |D_M(z)|\le2^{-4M^2-6M-1}+2^{-4M^2-2M+1}
           \le2^{-4M^2-2M+2}.
\]
On $|z|=1$, real coefficients give $B_M(z)=\overline{A_M(z)}$, which proves
exact outer-circle equality. On $|z|=1/2$ one has $\overline z=b/z$, so
\[
 \begin{split}
 \overline{A_M(z)}&=b^{-2}z^4B_M(z)+D_M(1/z),\\
 \overline{B_M(z)}&=b^{-2}z^4\{A_M(z)-D_M(z/b)\}.
 \end{split}
\]
Multiplication and subtraction eliminate the common term. As
$|b^{-2}z^4|=1$, $|A_M(z)|,|B_M(z)|<4$, and both $1/z$ and $z/b$ have
modulus two, this bounds the unnormalized intensity difference by
$2^{-4M^2-2M+5}$. The squared normalization factor on this circle is
$2^{-4L-4}$, proving the stated inner-circle estimate. It is uniform in the
angle; no angular sampling is involved.

At $x=3/4$, positivity of coefficients and three terms of $B_M$ give
\[
 B_M(x)\ge\frac43+\frac{27}{256}+\frac{16}{243}>\frac32.
\]
For a checkable infinite-sum upper bound on $A_M(x)$ put
$T_n=b^{2n^2+n}x^{4n+1}$. Consecutive ratios bound both omitted tails and give
\[
 A_M(x)\le A_\infty(x)
 \le\sum_{n=-2}^{1}T_n+
       \frac{T_2}{1-b^{11}x^4}+\frac{T_{-3}}{1-b^{13}x^{-4}}
 =\frac{86380343555303460998075}{64074215137581349797888}
 <\frac{27}{20}.
\]
Thus $B_M(x)^2-A_M(x)^2>2/5$. Restoring the squared normalization proves
$|H_M(x)|\ge (3/8)^{2L}/40$. The logarithm of the ratio of this lower bound
to any fixed positive power of the measurement upper bound tends to positive
infinity quadratically in $M$, proving the absence of a uniform H\"older bound.
\end{proof}

The theorem concerns absolute intensity errors, not phase error at a fixed
nonvanishing reference function. Its examples do not keep the measured
intensities bounded below. Additional nonvanishing or finite-dimensional
conditioning assumptions are not ruled out as routes to better stability.

\begin{corollary}[A lower bound on any uniform modulus]
Let $\Omega(\delta)$ be the supremum of the target intensity error at $3/4$
over pairs bounded by one on $D_2$ with measurement error at most $\delta$
on the two circles. Then $\Omega(\delta)\to0$ as $\delta\downarrow0$, but,
for some positive constants $c,C$ and all sufficiently small $\delta$,
\[
 \Omega(\delta)\ge c\exp\{-C\sqrt{\log(1/\delta)}\}.
\]
No matching upper rate is asserted.
\end{corollary}
\begin{proof}
If the limit were not zero, Montel's theorem would provide locally uniformly
convergent subsequences of a violating sequence of bounded pairs. The limits
have equal intensity on both entire circles, hence are constant-phase
multiples by Perez's theorem or Theorem~\ref{thm:holomorphic-identify}.
Their target intensity errors therefore tend to zero, a contradiction.
For the lower bound take
$M=\lceil\tfrac12\sqrt{\log_2(1/\delta)}\rceil$ in the theorem. Its
measurement upper bound is at most $2^{-4M^2}\le\delta$, while its target
lower bound has logarithm bounded below by a negative constant times
$1+\sqrt{\log(1/\delta)}$.
\end{proof}

Severe ill-posedness of infinite-dimensional phase retrieval is not a new
phenomenon: Alaifari and Grohs \cite{AG2021} prove quadratic-exponential
instability in a Gabor setting. Our construction concerns a different,
explicit observation geometry, namely two complete circles inside a fixed
disk and an interior target. The elementary recurrence
\eqref{eq:theta-truncation-defect} and exact rational bounds are the input
here, not a claim to have discovered general phase-retrieval instability.
The accompanying finite verifier re-expands the full angular intensity
polynomial and checks its coefficient $\ell^1$ bound independently of this
recurrence. Neither the inverse-problem result nor its numerical illustration
is a BGS spectral-statistics theorem.
